\documentclass[12pt]{article}

\usepackage[a4paper,margin=2.5cm]{geometry}
\usepackage[T1]{fontenc}
\usepackage{lmodern}
\usepackage{microtype}
\usepackage{hyperref}
\usepackage{amssymb}
\usepackage{makecell}
\usepackage{adjustbox}
\usepackage{multirow}
\usepackage{enumitem}
\usepackage{tabularx}
\usepackage{array}
\usepackage{longtable}
\usepackage{xcolor}
\definecolor{topicpurple}{HTML}{7030A0}
\definecolor{Purple}{HTML}{7030A0}
\newcommand{\purple}[1]{\textcolor{topicpurple}{\textit{#1}}}

\newcommand{\field}[1]{\noindent #1\par\vspace{1.2em}}
\newcommand{\checkbox}{\(\square\)}
\newcommand{\checked}{\(\checkmark\)}

\title{Topic Definition}
\date{}
\author{Alex Gu \\ 48028903}
\begin{document}
\maketitle
\subsection*{Administrative details}
\field{Student name: Alex Gu}
\field{Student number: 48028903}
\field{Program enrolment: Bachelor of Computer Science (Honours)}
\field{Course enrolment: REIT4881}
\field{Number of units: 8}
\field{Expected workload for your enrolment: 20 hours/week}

\subsection*{Supervision details}
\field{Primary Supervisor: Ruihong Qiu}
\field{Other Supervisors: Zi (Helen) Huang}
\field{Have you met your supervisor: Yes - in person}
\field{Do you have a regular meeting scheduled: Weekly}

\subsection*{Project title and topic}
\field{\textbf{Topic title:}}
\field{\indent{Privacy-Preserving Federated Learning for Distributed Skin Disease Classification}}
\field{\textbf{Topic description:}}
\indent{Federated learning enables distributed model training across multiple nodes ((independent devices or servers that each hold local data)), where each node trains locally on its own data. However, prior work has shown that standard federated averaging can struggle with non-identically distributed (non-IID) data and that shared model updates can still leak sensitive patient information through attacks.

This project investigates federated learning for multi-class skin disease classification. An image dataset (~20k images) covering 30 skin disease classes will be partitioned across multiple nodes, simulating a distributed healthcare environment. This project will implement federated learning methods, specifically FedAvg \cite{mcmahanCommunicationEfficientLearningDeep2023} and FedProx \cite{liFederatedOptimizationHeterogeneous2020}, and compare their performance against a centralised training baseline. The project will evaluate classification accuracy, communication efficiency, and model behaviour under non-IID data distributions across nodes.}

\subsection*{Societal relevance and research aim}
\field{\textbf{Why is your research topic important to society?}}
\field{\indent{Skin diseases affect millions of people worldwide and AI has shown strong potential in supporting diagnoses through image classification. However, building reliable models requires large, diverse datasets that are often distributed across institutions. Privacy regulations such as HIPAA and GDPR prevent centralised data aggregation, limiting the development of models that generalise across patient populations. Privacy-preserving federated learning enables collaboration without exposing sensitive data, potentially improving accuracy while maintaining patient trust and regulatory compliance.}}
\field{\textbf{Describe your research question or aim:}}
\indent{The aim of this project is to evaluate federated learning approaches for multi-class skin classification in a simulated distributed healthcare environment. Training medical image classifiers typically requires aggregating data on a single centralised server, which is often infeasible due to privacy concerns. Federated learning keeps data local and shares only model updates, but standard federated averaging lacks convergence guarantees on non-IID data and may still leak patient information.}\par
\indent{This project addresses the following questions: How do FedAvg and FedProx compare on a 30-class skin disease task when data is heterogeneously distributed across nodes? How do these approaches compare against a centralised baseline in terms of classification accuracy and communication efficiency? If time permits, improvements such as stronger architectures or communication optimisations may be explored.}

\subsection*{Referencing tools and style}
\field{What reference manager are you using: Zotero with Better BibTeX plugin, exporting to BibTeX (.bib) for use with LaTeX}
\field{What reference style are you going to use: IEEE}
\field{Why: IEEE is the standard citation style for computer science and engineering research. It is also the format used by most papers within this domain.}

\newpage

\subsection*{Reference list task}
This project investigates federated learning for distributed skin disease classification. McMahan et al. \cite{mcmahanCommunicationEfficientLearningDeep2023} introduced Federated Averaging (FedAvg), a foundational algorithm that enables collaborative training across decentralised devices by sharing model updates with a central server rather than raw data. Li et al. \cite{liFederatedLearningChallenges2020} surveyed the challenges surrounding federated learning, identifying key issues such as expensive communication, system heterogeneity, data heterogeneity from non-IID data and potential security risks associated with shared model updates. To address the convergence issue of FedAvg under heterogeneity, Li et al. \cite{liFederatedOptimizationHeterogeneous2020} proposed FedProx, which introduces a proximal term that prevents the decentralised models from deviating too far from the global model. This provides the convergence guarantee that FedAvg lacks. While FedAvg keeps raw data local, Kaissis et al. \cite{kaissisSecurePrivacypreservingFederated2020} highlighted that model updates could still leak patient information through attacks such as model inversion (rebuilding data from model weights) and membership inference (identifying if data is present within a set of training data). They argued that federated learning alone is not sufficient to meet the privacy mechanisms essential for medical imaging under regulations such as the Health Insurance Portability and Accountability Act (HIPAA) and the General Data Protection Regulation (GDPR). Bonawitz et al. \cite{bonawitzPracticalSecureAggregation2017} addressed this by designing a protocol that enables the server to calculate the sum of client updates without observing any individual, maintaining privacy even when clients drop out. Together, these works inform the design and evaluation of the federated learning framework proposed in this project.

\bibliographystyle{ieeetr}
\bibliography{references}

\subsection*{Use of AI statement}
\subsubsection*{}
\noindent \textit{I acknowledge the use of generative AI tools in completing this assessment. Details of which tools were used and how they were used are provided in the table below, along with appropriate in-text and full references. I take responsibility for critically evaluating and integrating the AI-generated content, and ensuring it adheres to academic integrity standards.}

\subsubsection*{Have you used AI to}
\begin{itemize}[leftmargin=1.5em]
  \item Explore your topic? Yes -- discussed federated learning and privacy concepts with Claude to build understanding.
  \item Define your research question? No -- research question was defined through meetings with supervisors.
  \item Assist with your literature review? Yes -- used Claude to generate notes, verified against originals.
  \item Assist with formatting the answers presented here? Yes -- used Claude for LaTeX troubleshooting
  \item Other aspects of your work? No
\end{itemize}

\subsubsection*{}


% --- the table ---
\begin{tabular}{|c|c|c|c|c|c|c|c|c|c|}
\hline
\multirow{2}{*}{\parbox{2.5cm}{\vspace*{\fill}\centering AI Model and date\vspace*{0pt}}}
 & \multicolumn{9}{c|}{\parbox{11cm}{\centering Table 1. AI tool usage (tick all that apply)}} \\ \cline{2-10}
 & \adjustbox{rotate=90}{Language Translation}
 & \adjustbox{rotate=90}{Grammar / style / spelling}
 & \adjustbox{rotate=90}{Topic exploration}
 & \adjustbox{rotate=90}{Research question}
 & \adjustbox{rotate=90}{Literature review}
 & \adjustbox{rotate=90}{Content Creation visual (formatting)}
 & \adjustbox{rotate=90}{Content Creation (e.g. code, text)}
 & \adjustbox{rotate=90}{Feedback}
 & \adjustbox{rotate=90}{Other (provide details)} \\ \hline
\makecell{Claude Opus\\Week 2-3}
 & & & \checkmark & & \checkmark & \checkmark & & & \\ \hline
 & & & & & & & & & \\ \hline
 & & & & & & & & & \\ \hline
\end{tabular}


\renewcommand{\arraystretch}{1.2}
\end{document}
